\documentclass[11pt]{article}
\usepackage[T1]{fontenc}
\usepackage{lmodern}
\usepackage[margin=1.05in]{geometry}
\usepackage{amsmath,amssymb,amsthm,mathtools}
\usepackage{microtype}
\usepackage{aliascnt}
\usepackage[colorlinks=true,linkcolor=blue,citecolor=blue,urlcolor=blue]{hyperref}
\usepackage[nameinlink,noabbrev]{cleveref}
\newtheorem{theorem}{Theorem}[section]
\newaliascnt{proposition}{theorem}
\newtheorem{proposition}[proposition]{Proposition}
\aliascntresetthe{proposition}
\newaliascnt{lemma}{theorem}
\newtheorem{lemma}[lemma]{Lemma}
\aliascntresetthe{lemma}
\theoremstyle{remark}
\newaliascnt{remark}{theorem}
\newtheorem{remark}[remark]{Remark}
\aliascntresetthe{remark}
\numberwithin{equation}{section}
\newcommand{\N}{\mathbb N}
\newcommand{\E}{\mathbb E}
\newcommand{\Penv}{\mathbb P}
\title{Strong disorder and critical-point shift\\for a Gaussian pinning model}
\author{Kaiwen Jin\thanks{Email: \texttt{yui.ui.siki@gmail.com}.}}
\date{October 15, 2026}
\begin{document}
\maketitle
\begin{abstract}
We study renewal pinning with independent standard Gaussian disorder and gap
law $K(n)=n^{-5/4}/\zeta(5/4)$. For every inverse temperature $\beta>0$,
we prove that almost-sure extinction of the normalized free partition
function at the annealed field $-\beta^2/2$ implies a strictly larger quenched
critical field. The decreasing renewal hazard allows us to control future
intersections through the polymer's age distribution. Under strong disorder,
this gives a $1/t$ lower tail for the maximum constrained partition function.
Testing favorable stretches along a planted renewal, while retaining the
final renewal endpoint, then gives summable fractional moments and a positive
change of field with zero free energy. A quantitative estimate of Alexander
and Berger excludes extinction at and below the second-moment boundary.
\end{abstract}

\section{Introduction}\label{sec:intro}

In a renewal pinning model, the contact times of a polymer with a defect
line form a renewal process. Random rewards at the contacts compete with the
entropy of long excursions. The critical field marks the onset of positive
free energy, while the normalized partition function at the annealed field
provides a martingale test for disorder. These two descriptions of the
transition need not coincide for a general disordered system.

We consider the following fixed renewal law. Let
$\tau=\{0=\tau_0<\tau_1<\cdots\}$ have iid gaps with
\begin{equation}\label{eq:K}
 K(n):=P_\tau(\tau_1=n)=\frac{n^{-5/4}}{\zeta(5/4)},\qquad n\ge1,
 \quad \zeta(5/4)=\sum_{m\ge1}m^{-5/4}.
\end{equation}
Let $(\omega_n)_{n\ge1}$ be iid standard real Gaussians, independent of
$\tau$, and write $\Penv,\E$ for their law and expectation. For
$\beta>0$ and $h\in\mathbb R$, define the free partition function
\begin{equation}\label{eq:model}
 Z_{N,h}^{\beta,\omega}=E_\tau\exp\!\left(
        \sum_{n=1}^N(\beta\omega_n+h)1_{\{n\in\tau\}}\right),
 \qquad
 F(\beta,h)=\lim_{N\to\infty}\frac1N\E\log Z_{N,h}^{\beta,\omega},
\end{equation}
and $h_c(\beta)=\inf\{h:F(\beta,h)>0\}$. The endpoint at $N$ is not
constrained to be a contact. The annealed critical field is
$h_a(\beta)=-\beta^2/2$. The process
$W_N=Z_{N,h_a}^{\beta,\omega}$ is a nonnegative mean-one martingale,
as follows by averaging the next Gaussian contact weight. Write $W_\infty$
for its almost-sure limit. We call
$\Penv(W_\infty=0)=1$ strong disorder.

The martingale approach establishes disorder irrelevance in a small-disorder
range for renewal exponents below $1/2$; see Lacoin~\cite{Lacoin}.
Fractional-moment estimates give a strict critical-point shift for
sufficiently large Gaussian disorder~\cite{Toninelli}. Neither statement
alone identifies what happens at every inverse temperature for which the
martingale becomes extinct. Berger~\cite[Conjecture~4.9]{Berger} asks whether
strong disorder forces a strict critical-point shift for pinning laws with
nonzero renewal exponent. Our result treats the exact law \eqref{eq:K}.

\begin{theorem}\label{thm:main}
For the renewal law \eqref{eq:K} and independent standard Gaussian disorder,
for every $\beta>0$,
\[
 \Penv(W_\infty=0)=1
 \quad\Longrightarrow\quad h_c(\beta)>-\beta^2/2.
\]
More precisely, under this premise there is $\Delta_\beta>0$ such that
$F(\beta,-\beta^2/2+\Delta_\beta)=0$.
\end{theorem}

Thus extinction at the annealed field is accompanied by a genuine interval
of zero free energy above that field. The positive interval may depend on
$\beta$. The theorem uses the exact normalized power law; it makes no claim
for arbitrary slowly varying modifications of its tail.

The proof follows the strategy of finding large partition weights and
detecting them under size bias, as used for bulk directed polymers by Junk
and Lacoin~\cite{JL}. In renewal pinning, however, the disorder is collected
only at contacts, and long excursions have polynomial probabilities. To
control where the polymer can collect disorder, we use its age since the
last contact. The decreasing hazard in \eqref{eq:K} makes the squared future
contact probabilities uniformly small when the next contact probability is
small. Above the second-moment boundary, their conditional drift forces a
large constrained weight before a high martingale level is reached.

Such weights can be found in logarithmic windows. Under a Gaussian size
bias, the environment contains a planted renewal. Tests begun at its renewal
times have iid failure-and-length pairs. Keeping the final planted endpoint
in this calculation preserves the power $K(N)$; this is what makes the
resulting fractional moments summable and permits a strictly positive change
of field. The second-moment boundary itself is handled separately, using
the quantitative coarse-graining estimate for one renewal pinned on another
in Alexander and Berger~\cite{AB}.

Section~\ref{sec:age} develops the age estimate, and
Section~\ref{sec:tail} proves the large-weight bound.
Section~\ref{sec:field} obtains the positive field change.
Section~\ref{sec:boundary} treats the remaining boundary and completes the
proof. Constants may depend on the fixed $\beta$ and on explicitly fixed
auxiliary exponents. We suppress $\beta,\omega$ from partition functions
when no ambiguity is possible.

\section{The age process and future intersections}\label{sec:age}

Throughout the proof, $s=5/4$, $\alpha=s-1=1/4$ and
$c=\zeta(s)^{-1}$, so $K(n)=cn^{-s}$. Write
\[
 \eta_n=e^{\beta\omega_n-\beta^2/2},\qquad
 \xi_n=\eta_n-1,\qquad \chi=\E\xi_n^2=e^{\beta^2}-1,
 \qquad \mathcal F_n=\sigma(\omega_1,\ldots,\omega_n).
\]
All absolute moments of $\xi_n$ are finite and $\xi_n>-1$.
Let $C_n$ be the partition function constrained to have a contact at $n$,
at field $h_a=-\beta^2/2$. Then
\begin{equation}\label{eq:renewal-recursion}
 C_0=1,\qquad C_n=\eta_n\sum_{j=0}^{n-1}K(n-j)C_j,
 \qquad W_N=\sum_{j=0}^N C_j\overline K(N-j),
 \quad \overline K(a)=\sum_{m>a}K(m).
\end{equation}
In particular, $C_n\le W_n$, $\E C_n=u(n):=P_\tau(n\in\tau)$,
and $u(0)=\overline K(0)=1$. The renewal theorem
\cite[Theorem~3.16]{Berger} gives $u(n)\asymp n^{-3/4}$. Hence
\begin{equation}\label{eq:R}
 0<R:=\sum_{n\ge1}u(n)^2<\infty,\qquad
 \beta_2^2:=\log(1+1/R).
\end{equation}
The strict inequality $\beta>\beta_2$ is equivalent to $\chi R>1$.
We first prove the critical-point shift in that range under strong disorder.
Section~\ref{sec:boundary} will exclude strong disorder at and below $\beta_2$.

The age of a renewal is the time elapsed since its last contact. Its
transition matrix $T$ on $\N_0$ resets age $a$ to zero with probability
\[
 r(a)=\frac{K(a+1)}{\overline K(a)}
\]
and otherwise moves to $a+1$. The following property of the exact law will
control contacts far into the future.

\begin{lemma}\label{lem:age}
The hazard $r(a)$ decreases to zero and satisfies $r(a)\le\alpha/(a+1)$.
For $u_a(k):=T^k(a,0)$, $k\ge1$, one has
\[
 0\le u_a(k)\le u(k),\qquad
 \sum_{k\ge1}u_a(k)^2\longrightarrow0\quad(a\to\infty).
\]
\end{lemma}
\begin{proof}
With $m=a+1$,
\[
 r(a)^{-1}=\sum_{j\ge0}\left(\frac{m}{m+j}\right)^s.
\]
Each summand increases with $m$. Also
$\sum_{n=m}^\infty n^{-s}\ge\int_m^\infty x^{-s}\,dx$,
which gives the asserted upper bound on $r$.
Couple chains begun at ages $a$ and zero with the same uniform variables.
The larger age remains larger: whenever it resets, the smaller age resets
as well, and the chains then agree. Thus $u_a(k)\le u_0(k)=u(k)$.
For fixed $k$, a reset before time $k$ has probability at most $kr(a)$,
so $u_a(k)\to0$. Dominated summation using \eqref{eq:R} proves the last claim.
\end{proof}

For a probability mass $\mu$ on ages, put
\begin{equation}\label{eq:potential}
 f_\mu(k)=(\mu T^k)(0),\quad k\ge1,\qquad
 U(\mu)=\sum_{k\ge1}f_\mu(k)^2.
\end{equation}
This quantity measures the expected future overlap of two independent
renewals begun with age law $\mu$. We use the same linear definition of
$f_\mu$ for differences of the probability masses appearing below.
By \cref{lem:age}, $\|f_{\delta_a}\|_2\le\sqrt R$, so these sums are
well defined and $0\le U(\mu)\le R$.

\begin{lemma}\label{lem:compactness}
If $q=(\mu T)(0)$ tends to zero, then $U(\mu)$ and $U(\mu T)$ tend
to zero uniformly over all probability masses $\mu$. Moreover,
\begin{equation}\label{eq:potential-shift}
 U(\mu T)=U(\mu)-q^2.
\end{equation}
\end{lemma}
\begin{proof}
For every integer $A\ge0$, monotonicity of the hazard gives
\[
 \mu([0,A])\le q/r(A),\qquad
 \|f_\mu\|_2\le\sqrt R\,q/r(A)
                  +\sup_{a>A}\|f_{\delta_a}\|_2.
\]
Choose $A$ large by \cref{lem:age}, and then choose $q$ small. This proves
uniform convergence. Finally $f_{\mu T}(k)=f_\mu(k+1)$, giving
\eqref{eq:potential-shift} and the assertion for $\mu T$.
\end{proof}

Under the polymer measure at time $n$, the age law is
\[
 \mu_n(a)=\frac{C_{n-a}\overline K(a)}{W_n},\qquad 0\le a\le n.
\]
Given the past, let $v=\mu_{n-1}T$, $q_n=v(0)$ and
$z=q_n\xi_n/(1+q_n\xi_n)$. Prediction by $T$, followed by reweighting the
mass at zero by $\eta_n$, gives the exact identities
\begin{equation}\label{eq:age-update}
 \frac{W_n}{W_{n-1}}=1+q_n\xi_n,\qquad
 \mu_n=v+z(\delta_0-v).
\end{equation}
Here $0<q_n\le1/4$ and $q_n$ is $\mathcal F_{n-1}$-measurable.
Let $\Delta U_n=U(\mu_n)-U(\mu_{n-1})$. Expanding the square in
\eqref{eq:potential} gives
\begin{equation}\label{eq:deltaU}
 \Delta U_n=-q_n^2+2zH+z^2J,
 \quad H=\langle f_v,f_{\delta_0-v}\rangle,
 \quad J=\|f_{\delta_0-v}\|_2^2.
\end{equation}

\begin{lemma}\label{lem:exponential}
Suppose $\beta>\beta_2$. Set $\gamma=(\chi R-1)/2>0$. For every
$\lambda>0$ there is $\delta>0$ such that, whenever $q_n\le\delta$,
\begin{equation}\label{eq:exponential}
 \E[e^{-\lambda\Delta U_n}\mid\mathcal F_{n-1}]
 \le e^{-a q_n^2},\qquad a=\lambda\gamma/2.
\end{equation}
Thus $a$ can be made arbitrarily large, with a corresponding smaller $\delta$.
\end{lemma}
\begin{proof}
For $q=q_n\le1/4$, the denominator $1+q\xi_n$ is at least $3/4$.
Conditionally on the past,
\[
 |z|\le\tfrac43q|\xi_n|,\quad |z|\le1,\quad
 \E z=-q^2\E\frac{\xi_n^2}{1+q\xi_n}=O(q^2),\quad
 \E z^2=\chi q^2+O(q^3),\quad \E|z|^4=O(q^4).
\]
The error in the second moment follows by bounding
$|(1+x)^{-2}-1|\le C|x|$ for $x\ge-1/4$.
By \cref{lem:compactness}, uniformly as $q\downarrow0$,
\[
 |H|\le\sqrt{U(v)}(\sqrt R+\sqrt{U(v)})\longrightarrow0,
 \qquad J\longrightarrow R,\qquad 0\le J\le4R.
\]
Equation~\eqref{eq:deltaU} therefore implies
\begin{equation}\label{eq:drift}
 \frac{\E[\Delta U_n\mid\mathcal F_{n-1}]}{q^2}\longrightarrow\chi R-1,
 \qquad
 \frac{\E[(\Delta U_n)^2\mid\mathcal F_{n-1}]}{q^2}\longrightarrow0.
\end{equation}
For the latter assertion use
$(\Delta U_n)^2\le3q^4+12H^2z^2+3J^2z^4$.
Also $\Delta U_n\ge-q^2-2|H|\ge-b(\delta)$ for $q\le\delta$,
where $b(\delta)\to0$.
For $x\ge-b$ the exponential remainder satisfies
\[
 e^{-\lambda x}\le1-\lambda x+\tfrac12\lambda^2e^{\lambda b}x^2.
\]
Choose $\delta$ so the first conditional mean in \eqref{eq:drift} is at
least $\gamma q^2$ and the remainder contributes at most
$\lambda\gamma q^2/2$. Then the conditional expectation is at most
$1-aq^2\le e^{-aq^2}$.
\end{proof}

In particular, before a step with $q_n>\delta$,
\begin{equation}\label{eq:potential-martingale}
 \exp\!\left\{-\lambda U(\mu_n)+a\sum_{i\le n}q_i^2\right\}
\end{equation}
is a nonnegative supermartingale. This stopping is legitimate before the
exceptional step because $q_n$ is known at time $n-1$.

\section{Large constrained weights}\label{sec:tail}

In this section assume strong disorder, $\Penv(W_\infty=0)=1$, and
$\beta>\beta_2$. The aim is a $1/t$ lower bound for the tail of
$\sup_n C_n$. This will produce detectable favorable stretches whose
length is only logarithmic in their weight.

\begin{lemma}\label{lem:overshoot}
Let $\tau_x=\inf\{n\ge0:W_n\ge x\}$, with $\inf\varnothing=\infty$.
For every fixed $p\ge1$ there is $M_p<\infty$ such that, for $x\ge1$,
\begin{equation}\label{eq:hitting}
 \frac1{M_1x}\le\Penv(\tau_x<\infty)\le\frac1x,\qquad
 \E[W_{\tau_x}^p;\tau_x<\infty]\le M_p x^{p-1}.
\end{equation}
\end{lemma}
\begin{proof}
First, lognormal relative overshoots have uniformly bounded moments:
\begin{equation}\label{eq:relative-overshoot}
 \sup_{t\ge1}\E[(\eta_1/t)^p\mid\eta_1\ge t]<\infty.
\end{equation}
Indeed, with $b=(\log t+\beta^2/2)/\beta\ge\beta/2$, the conditional
expectation equals
\[
 \frac{\int_0^\infty e^{p\beta y-by-y^2/2}\,dy}
      {\int_0^\infty e^{-by-y^2/2}\,dy}.
\]
For $b\ge\max(1,2p\beta)$ the numerator is at most $(b-p\beta)^{-1}$
and the denominator is at least $e^{-3/2}/b$, by restricting to $[0,1/b]$.
On the remaining bounded interval, the numerator is uniformly finite and
the denominator is bounded away from zero.

At a first crossing of $x>1$, conditional on the past, the condition is
$\eta_n\ge t=(x/W_{n-1}-1+q_n)/q_n>1$. By \eqref{eq:age-update},
\[
 \frac{W_n}{x}=\frac{1-q_n+q_n\eta_n}{1-q_n+q_nt}\le\eta_n/t.
\]
Apply \eqref{eq:relative-overshoot} conditionally and sum over the disjoint
first-crossing times. This gives
\[
 \E[W_{\tau_x}^p;\tau_x<\infty]\le M_p x^p\Penv(\tau_x<\infty).
\]
The nonnegative martingale maximal inequality gives the upper probability
bound in \eqref{eq:hitting}. Bounded optional stopping yields
\[
 1=\E[W_{\tau_x};\tau_x\le n]+\E[W_n;\tau_x>n].
\]
The second term tends to zero by strong disorder and domination by $x$.
Thus $\E[W_{\tau_x};\tau_x<\infty]=1$, proving the lower probability
bound and then the moment bound. For $x=1$, $\tau_1=0$; enlarge the constants
to include this case.
\end{proof}

\begin{proposition}\label{prop:highC}
There is $c_\beta>0$ such that
\begin{equation}\label{eq:highC}
 \Penv\left(\sup_{n\ge0}C_n\ge t\right)\ge c_\beta/t,
 \qquad t\ge1.
\end{equation}
\end{proposition}
\begin{proof}
Expanding the fourth power in \eqref{eq:age-update}, using $\E\xi_n=0$,
gives a constant $D>0$ with
\[
 \E(1+q\xi_n)^4\le e^{Dq^2}\quad(0\le q\le1/4).
\]
Hence $W_n^4\exp(-D\sum_{i\le n}q_i^2)$ is a nonnegative
supermartingale. Choose $\varepsilon_0>0$ so $4-D\varepsilon_0\ge3$.
Choose $\lambda$, then $\delta$, in \cref{lem:exponential} so that
$a\varepsilon_0\ge3$.

Let $B>1$ be a level multiplier, to be fixed below. On the event that
$\tau_{Bx}<\infty$, write $A=\sum_{n=\tau_x+1}^{\tau_{Bx}}q_n^2$.
The sum is empty if a single overshoot crosses both levels. Stopping the
fourth-moment supermartingale after $\tau_x$ and using
\cref{lem:overshoot} gives
\begin{equation}\label{eq:small-budget}
 \Penv(\tau_{Bx}<\infty,A\le\varepsilon_0\log B)
 \le M_4x^{-1}B^{-4+D\varepsilon_0}\le M_4x^{-1}B^{-3}.
\end{equation}
At the terminal crossing, its fourth-power factor is at least $(Bx)^4$
and its compensator at least $B^{-D\varepsilon_0}$, whereas the initial
expected fourth power is at most $M_4x^3$.

Similarly \eqref{eq:potential-martingale} and $0\le U\le R$ imply
\begin{equation}\label{eq:large-budget}
 \Penv\left(\tau_{Bx}<\infty,A\ge\varepsilon_0\log B,
       q_n\le\delta\ \forall\,\tau_x<n\le\tau_{Bx}\right)
 \le x^{-1}e^{\lambda R}B^{-a\varepsilon_0}
 \le x^{-1}e^{\lambda R}B^{-3}.
\end{equation}
To obtain this bound, condition at $\tau_x$ and stop before the first
$q_n>\delta$; the terminal exponential on the indicated event is at least
$e^{-\lambda R}B^{a\varepsilon_0}$ times its initial value. Multiply the
conditional probability bound by $\Penv(\tau_x<\infty)\le1/x$.
Both stopping estimates follow first at bounded horizons and then by Fatou.

Let $\sigma=\inf\{n\ge\tau_x:W_n\le x/B\}$. Conditional on a dip
before the upper hit, the probability of a later passage to $Bx$ is at most
$W_\sigma/(Bx)\le B^{-2}$. Therefore
\begin{equation}\label{eq:no-dip}
 \Penv(\sigma<\tau_{Bx}<\infty)\le x^{-1}B^{-2}.
\end{equation}
By \eqref{eq:hitting}, $\Penv(\tau_{Bx}<\infty)\ge(M_1Bx)^{-1}$.
Fix $B$ sufficiently large that subtracting
\eqref{eq:small-budget}--\eqref{eq:no-dip} leaves a probability at least
$c_*/x>0$. On the remaining event some $\tau_x<n\le\tau_{Bx}$ has
$q_n>\delta$, and $W_{n-1}\ge x/B$. Thus
\[
 \max_{i<n}C_i\ge\sum_{i<n}K(n-i)C_i=q_nW_{n-1}\ge\delta x/B,
\]
where the first inequality uses $\sum_{i<n}K(n-i)\le1$.
Taking $x=Bt/\delta$ proves \eqref{eq:highC}.
\end{proof}

\begin{lemma}\label{lem:log-window}
For every $\varepsilon>0$ there is $D_\varepsilon<\infty$ such that, for
all sufficiently large $t$, with $L=\lceil D_\varepsilon\log t\rceil$,
\begin{equation}\label{eq:log-window}
 E_{t,L}:=\{\max_{0\le n\le L}C_n\ge t\},\qquad
 \Penv(E_{t,L})\ge t^{-1-\varepsilon}.
\end{equation}
\end{lemma}
\begin{proof}
Constrained concatenation gives
$C_{n+m}(\omega)\ge C_n(\omega)C_m(\vartheta^n\omega)$, where
$(\vartheta^n\omega)_j=\omega_{n+j}$. Choose a large fixed $v>1$ with
$c_\beta/(2v)\ge v^{-1-\varepsilon/2}$, then a finite $M$ with
$\Penv(\max_{n\le M}C_n\ge v)\ge c_\beta/(2v)$, by
\cref{prop:highC}. Repeat this success at its first hitting time, using fresh
future Gaussian coordinates. The probability of $k$ successive successes
is at least $v^{-k(1+\varepsilon/2)}$, their total time is at most $kM$,
and concatenation gives a value at least $v^k$.
Take $k=\lceil\log_v t\rceil$. For large $t$ the fixed extra factor is
absorbed in $t^{-\varepsilon/2}$, and $kM\le D_\varepsilon\log t$ after
enlarging $D_\varepsilon$.
\end{proof}

\section{A positive field shift}\label{sec:field}

We use a change of environment law that makes favorable stretches more
likely. Sample a renewal $\sigma$ with law $P_\tau$ and independent standard
Gaussians $(G_n)$, and set
\[
 \omega_n=G_n+\beta1_{\{n\in\sigma\}}.
\]
Denote the resulting joint law by $Q$. Completing the square in the Gaussian
density gives, for every nonnegative measurable function $f$ of the first
$N$ environment coordinates,
\begin{equation}\label{eq:planting}
 E_Qf(\omega)=\E[W_Nf(\omega)],\qquad
 E_Q[1_{\{N\in\sigma\}}f(\omega)]=\E[C_Nf(\omega)].
\end{equation}
These are finite-dimensional identities. For $E_{t,L}$ in
\cref{lem:log-window}, stop $W$ at the first constrained crossing of $t$,
truncated at $L$. On success its stopped value is at least $t$. Therefore
\begin{equation}\label{eq:planted-success}
 p:=Q(E_{t,L})=\E[W_L;E_{t,L}]
 \ge t\Penv(E_{t,L})\ge t^{-\varepsilon}.
\end{equation}
On the original environment law, the martingale maximal inequality and
$C_n\le W_n$ give $\Penv(E_{t,L})\le1/t$.

\begin{lemma}\label{lem:failed-tests}
Let $\mathcal A_N$ be the event that $E_{t,L}$ occurs in some window
starting at an integer in $[0,N-L]$. Then
\begin{equation}\label{eq:failed-tests}
 \Penv(\mathcal A_N)\le(N+1)/t,\qquad
 \E[C_N;\mathcal A_N^c]\le C L^{s+1}p^{-(s+2)}K(N)
 \quad(N\ge2L).
\end{equation}
Here $C$ depends only on the renewal law.
\end{lemma}
\begin{proof}
The first bound is a union bound over deterministic window starts.
Under $Q$, begin a test at the planted renewal zero, use the next $L$
coordinates, and start the next test at
$T_L=\min\{\sigma_j:\sigma_j\ge L\}$. Iterate this rule.
At each start, the future renewal gaps and the Gaussian coordinates strictly
after that start are fresh. The pairs consisting of the failure indicator
and the next regeneration length are consequently iid. Failure may depend
on that length; we retain this dependence in
\[
 g(n)=Q(E_{t,L}^c,T_L=n),\qquad n\ge L,\qquad
 \sum_n g(n)=1-p.
\]
The test uses no Gaussian coordinate after its next regeneration, and a
subsequent test excludes the coordinate at its own start. Thus even when
$T_L=L$, the two tests use disjoint noise. Independence follows by
conditioning on each possible renewal-gap index and physical regeneration
time and applying iid independence.

The last renewal before $L$ gives
\begin{equation}\label{eq:g-envelope}
 g(n)\le P_\sigma(T_L=n)
   =\sum_{j=0}^{L-1}u(j)K(n-j)\le B_LK(n),
 \qquad B_L\le C L^s.
\end{equation}
For $n\ge2L$ use $K(n-j)\le2^sK(n)$ and $\sum_{j<L}u(j)\le L$.
For $L\le n<2L$, the probability is at most one whereas
$K(n)\ge c(2L)^{-s}$.

For any nonnegative kernel $g$ of mass $1-p$ bounded by $BK$, one has
\begin{equation}\label{eq:defective-convolution}
 g^{*m}(n)\le Bm^{s+1}(1-p)^{m-1}K(n),\qquad
 \sum_{m\ge1}g^{*m}(n)\le C_s Bp^{-(s+2)}K(n).
\end{equation}
Indeed one of the $m$ summands of a composition of $n$ is at least $n/m$.
Its factor is at most $Bm^sK(n)$. Sum the other factors using their total
mass and then sum over the $m$ possible positions. Finally
$\sum_{m\ge1}m^{s+1}(1-p)^{m-1}\le C_s p^{-(s+2)}$:
for $p\le1/2$ compare with a polynomial-exponential integral, and for
$p\ge1/2$ use a convergent geometric majorant. The case $p=1$, $g=0$,
is immediate.

On $\{N\in\sigma\}\cap\mathcal A_N^c$, every test starting at or
before $N-L$ fails. Its next regeneration is at most $N$, since $N$ itself
is a renewal at least $L$ beyond that start. The remaining distance $r$ after
the last test therefore satisfies $0\le r<L$. Set
$h_L(r)=u(r)1_{\{0\le r<L\}}$, with $g^{*0}$ the unit mass at zero.
By regeneration and \eqref{eq:planting},
\[
 \E[C_N;\mathcal A_N^c]=Q(N\in\sigma,\mathcal A_N^c)
 \le\sum_{m\ge0}(g^{*m}*h_L)(N).
\]
For $N\ge2L$ the $m=0$ term vanishes. In the other terms,
$N-r\ge N/2$, so $K(N-r)\le2^sK(N)$ and $\sum_rh_L(r)\le L$.
Apply \eqref{eq:defective-convolution} and \eqref{eq:g-envelope}.
\end{proof}

The factor $K(N)$ in \eqref{eq:failed-tests} supplies the summable power
needed below. It comes from retaining the planted endpoint $N\in\sigma$.

\begin{proposition}\label{prop:fractional-sum}
For some $\theta\in(4/5,1)$,
\begin{equation}\label{eq:fractional-sum}
 \sum_{N\ge0}\E C_N^\theta<\infty.
\end{equation}
\end{proposition}
\begin{proof}
Choose $\theta\in(1/s,1)$, then $b>1+1/(1-\theta)$ and
$\varepsilon>0$ so small that $\theta[s-b\varepsilon(s+2)]>1$.
Take $t=N^b$ and $L=\lceil D_\varepsilon\log t\rceil$. For all large
$N$, $N\ge2L$ and \eqref{eq:planted-success} gives $p\ge t^{-\varepsilon}$.
H\"older's inequality on $\mathcal A_N$ and its complement, together with
$\E C_N=u(N)\le1$, yields
\begin{equation}\label{eq:frac-bound}
 \E C_N^\theta\le
 \left(\frac{N+1}{t}\right)^{1-\theta}
 +\left[C L^{s+1}t^{\varepsilon(s+2)}K(N)\right]^\theta.
\end{equation}
The powers of $N$ are respectively $-(b-1)(1-\theta)<-1$ and
$-\theta[s-b\varepsilon(s+2)]<-1$, with only a fixed logarithmic factor
in the second term. Both series converge. The finitely many remaining
moments are finite by the lognormal moment formula.
\end{proof}

We finish this range of beta by turning \eqref{eq:fractional-sum} into an
actual positive change of field. Let $C_n^{(\Delta)}$ be the constrained
partition at $h_a+\Delta$, and put
$A_n(\Delta)=\E[(C_n^{(\Delta)})^\theta]$, $A_0(\Delta)=1$.

\begin{proposition}\label{prop:positive-field}
Under strong disorder and $\beta>\beta_2$, there exists $\Delta>0$ such
that $F(\beta,h_a+\Delta)=0$.
\end{proposition}
\begin{proof}
Fix an integer $k\ge1$. Decompose a path constrained at $N\ge k$ by its
last contact $j<k$ and first contact $i\ge k$. The factors
$C_j^{(\Delta)}$, $K(i-j)e^\Delta\eta_i$, and
$C_{N-i}^{(\Delta)}\circ\vartheta^i$ use disjoint environment coordinates.
Since $(\sum x_j)^\theta\le\sum x_j^\theta$ for nonnegative $x_j$,
\begin{align}
 A_N(\Delta)&\le\sum_{i=k}^N Q_i(\Delta)A_{N-i}(\Delta),
 \label{eq:finite-volume}\\
 Q_i(\Delta)&=\E\eta_1^\theta\,e^{\theta\Delta}
     \sum_{j<k}K(i-j)^\theta A_j(\Delta).
 \nonumber
\end{align}
By \cref{prop:fractional-sum} and $s\theta>1$,
\[
 \rho_k(0):=\sum_{i\ge k}Q_i(0)
 =\E\eta_1^\theta\sum_{j<k}A_j(0)
                         \sum_{r\ge k-j}K(r)^\theta\longrightarrow0.
\]
For each fixed $j$ the inner tail tends to zero and it is uniformly bounded
by $\sum_rK(r)^\theta<\infty$, so dominated summation applies.
Choose $k$ with $\rho_k(0)<1$. Since a path at $j$ has at most $j$
contacts, $A_j(\Delta)\le e^{\theta\Delta j}A_j(0)$, and hence
\[
 \rho_k(\Delta)\le e^{\theta\Delta k}\rho_k(0)<1
\]
for a sufficiently small, strictly positive $\Delta$.
Sum \eqref{eq:finite-volume} only up to a finite $M$ before rearranging:
\[
 \sum_{n\le M}A_n(\Delta)
 \le\frac{\sum_{j<k}A_j(\Delta)}{1-\rho_k(\Delta)}.
\]
It follows that $\sum_nA_n(\Delta)<\infty$.
The last-contact formula \eqref{eq:renewal-recursion}, with the additional
reward $e^\Delta$ at each contact, now gives
\[
 \E Z_{N,h_a+\Delta}^{\theta}
 \le\sum_{j\le N}A_j(\Delta)\overline K(N-j)^\theta
 \le\sum_{j\ge0}A_j(\Delta).
\]
Jensen's inequality bounds $\E\log Z_{N,h_a+\Delta}$ uniformly from above.
The no-contact paths give $\log Z_{N,h_a+\Delta}\ge\log\overline K(N)$,
whose ratio to $N$ tends to zero. Thus $F(\beta,h_a+\Delta)=0$.
\end{proof}

\section{The second-moment boundary}\label{sec:boundary}

We must still exclude strong disorder when $0<\beta\le\beta_2$.
We use a quantitative estimate for a renewal pinned on another renewal
from Alexander and Berger~\cite{AB}. Its fractional-moment form is important:
a statement only about a free-energy critical point would not give the
infinite partition function used here.

Let $\sigma$ and $\tau$ be independent renewals with gap law $K$. For
$\lambda>0$, define the partition constrained at the $j$th spine renewal by
\[
 V_j^\sigma(\lambda)=E_\tau\left[
 e^{\lambda\sum_{i=1}^j1_{\{\sigma_i\in\tau\}}};\ \sigma_j\in\tau\right],
 \qquad V_0^\sigma=1.
\]
The contact at zero is excluded. Put $z=e^\lambda-1$. Expansion over the
selected contacts gives
\begin{equation}\label{eq:chaos}
 \mathcal V_j^\sigma(\lambda):=(1-e^{-\lambda})V_j^\sigma(\lambda)
 =\sum_{m=1}^j z^m
   \sum_{0=i_0<i_1<\cdots<i_m=j}
       \prod_{a=1}^m u(\sigma_{i_a}-\sigma_{i_{a-1}}).
\end{equation}
The fixed factor is necessary because the forced endpoint has weight
$e^\lambda$ in $V$, but weight $z$ in the selected-contact expansion.
We retain it when using the chaos estimate in \cite[Section~4]{AB}, whose
displayed identification of the two constrained partitions omits that factor.

Tonelli's theorem gives
\[
 \sum_{j\ge1}E_\sigma u(\sigma_j)=\sum_{n\ge1}u(n)^2=R.
\]
Thus the reference kernel in \cite[(4.2)]{AB} is
$K^*(j)=R^{-1}E_\sigma u(\sigma_j)$, and the critical value of $z$ is
$z_0=1/R$, as also follows from \cite[Proposition~1.3]{AB}.
Write $\lambda_0=\log(1+1/R)=\beta_2^2$.

Here is the source estimate in the notation needed for the proof. For an
integer $L$, let $B_a=\{(a-1)L+1,\ldots,aL\}$. If
$I=\{1\le b_1<\cdots<b_\ell=n\}$, $b_0=0$, put
\begin{equation}\label{eq:restricted-chaos}
 Z_I^\sigma=\sum_{m=1}^{nL} R^{-m}
  \sum_{\substack{0=i_0<i_1<\cdots<i_m=nL\\
        \{\lceil i_a/L\rceil:1\le a\le m\}=I}}
       \prod_{a=1}^m u(\sigma_{i_a}-\sigma_{i_{a-1}}).
\end{equation}
This agrees with the restricted partition in \cite[(4.10)]{AB}: its
between-contact factor, written $K^*(b-a)w(\sigma,a,b)$ there, is exactly
$R^{-1}u(\sigma_b-\sigma_a)$.

\begin{proposition}[Alexander--Berger]\label{prop:AB}
Let $\zeta=3/4$. For every $\delta>0$ there is $L_0$ such that, for
every integer $L\ge L_0$, every $n\ge1$ and every such $I$,
\begin{equation}\label{eq:AB}
 E_\sigma[(Z_I^\sigma)^\zeta]
 \le\delta^\ell\prod_{a=1}^\ell(b_a-b_{a-1})^{-3/2}.
\end{equation}
\end{proposition}

This is \cite[(4.12)]{AB}, in the explicitly cited arXiv version.
Its hypotheses (1.5) hold with both slowly varying functions equal to $c$
and both exponents equal to $1/4$. In their notation,
$\alpha_*=(1-\alpha-\widetilde\alpha)/\widetilde\alpha=2>1/2$ and
$\zeta=(1+\alpha_*/4)/(1+\alpha_*/2)=3/4$; the exponent on block gaps
is $1+\alpha_*/4=3/2$. The arbitrary-small-$\delta$ assertion follows
there from Lemma~4.2 and (4.22), and is not restricted to a particular
mass-one normalization.

\begin{lemma}\label{lem:spine-sum}
There is $\lambda_1>\lambda_0$ such that
\begin{equation}\label{eq:spine-sum}
 \sum_{j\ge1}E_\sigma[V_j^\sigma(\lambda_1)^{3/4}]<\infty,
 \qquad \sum_{j\ge1}V_j^\sigma(\lambda_1)<\infty
 \quad P_\sigma\text{-a.s.}
\end{equation}
\end{lemma}
\begin{proof}
Choose $\delta>0$ so the kernel
$q(a)=e^\zeta\delta a^{-3/2}$ has mass $\rho<1$.
Fix an allowed $L$ in \cref{prop:AB} and put
$\lambda_1=\log(1+R^{-1}e^{1/L})>\lambda_0$.
In a term visiting $\ell$ blocks, at most $L\ell$ contacts are selected.
Changing $z_0$ to $z_0e^{1/L}$ therefore increases that term by at most
$e^\ell$. By the fractional-sum inequality, \eqref{eq:chaos} and
\eqref{eq:AB},
\[
 E_\sigma[V_{nL}^\sigma(\lambda_1)^\zeta]
 \le(1-e^{-\lambda_1})^{-\zeta}
            \sum_{\ell\ge1}q^{*\ell}(n).
\]
Summation over $n$ bounds the total by
$(1-e^{-\lambda_1})^{-\zeta}\rho/(1-\rho)<\infty$.

For an arbitrary spine index $j$, let $M=L(\lfloor j/L\rfloor+1)$.
A direct tau gap from $\sigma_j$ to $\sigma_M$ gives
\[
 V_M^\sigma(\lambda_1)\ge V_j^\sigma(\lambda_1)
                    e^{\lambda_1}K(\sigma_M-\sigma_j).
\]
The last factor uses only the next $M-j\in\{1,\ldots,L\}$ spine gaps,
independently of $V_j$. The minimum of the $L$ positive finite expectations
$E_\sigma[(e^{\lambda_1}K(\sigma_r))^\zeta]$, $1\le r\le L$, is positive.
Hence summability at multiples of $L$ implies summability at every $j$,
each multiple being used at most $L$ times. Tonelli gives a finite sum of
the $\zeta$ powers almost surely, and then a finite sum of $V_j$, since
eventually $V_j\le1$ and $V_j\le V_j^\zeta$.
\end{proof}

\begin{proposition}\label{prop:UI}
For every $0<\beta\le\beta_2$, the martingale $(W_N)$ is uniformly
integrable and $\E W_\infty=1$.
\end{proposition}
\begin{proof}
Expanding a finite product by its last selected factor and then increasing
the cutoff gives
\begin{equation}\label{eq:free-spine}
 V_\infty^{\sigma,\mathrm{free}}(\lambda_1)
 :=E_\tau e^{\lambda_1\sum_{j\ge1}1_{\{\sigma_j\in\tau\}}}
 =1+(1-e^{-\lambda_1})\sum_{j\ge1}V_j^\sigma(\lambda_1)<\infty
 \quad P_\sigma\text{-a.s.}
\end{equation}
Indeed $\prod_{i=1}^M e^{\lambda I_i}
=1+\sum_{j=1}^M(e^\lambda-1)I_j\prod_{i<j}e^{\lambda I_i}$;
on $I_j=1$, the missing endpoint reward converts the coefficient to
$1-e^{-\lambda}$. Monotone convergence and \cref{lem:spine-sum} apply.

Recall that under $Q$, $\omega=G+\beta1_\sigma$, where $G$ is iid standard
Gaussian and independent of $\sigma$. For every $\beta^2\le\lambda_1$, average
$W_N$ over $G$ as in \eqref{eq:planting}. The Gaussian moment formula gives
\[
 E_G[W_N(G+\beta1_\sigma)\mid\sigma]
   =E_\tau e^{\beta^2|\tau\cap\sigma\cap[1,N]|}
   \le V_\infty^{\sigma,\mathrm{free}}(\lambda_1).
\]
Conditional Markov's inequality and \eqref{eq:planting} imply, uniformly
in physical time $N$,
\[
 \E[W_N;W_N>A]=Q(W_N>A)
 \le E_\sigma\min\{1,V_\infty^{\sigma,\mathrm{free}}(\lambda_1)/A\}
 \longrightarrow0\quad(A\to\infty).
\]
The last limit is dominated convergence. Thus $W_N$ is uniformly
integrable and its almost-sure limit has mean one. Since
$\lambda_1>\beta_2^2$, this covers the asserted interval, including its
endpoint.
\end{proof}

\begin{proof}[Proof of \cref{thm:main}]
Fix $\beta>0$ with $\Penv(W_\infty=0)=1$. By \cref{prop:UI},
$\beta>\beta_2$. Proposition~\ref{prop:positive-field} then gives a
$\Delta>0$ with $F(\beta,h_a+\Delta)=0$. The free partition increases
with $h$, so $h_c(\beta)\ge h_a+\Delta>h_a$.
\end{proof}

\section*{AI assistance}
OpenAI Codex was used substantially in developing the mathematical argument,
checking source applications, and preparing the text. AI reader and
referee-style feedback was used in revising the text. Such feedback is not
independent human mathematical validation.

\begin{thebibliography}{9}
\raggedright
\bibitem{AB}
K.~S. Alexander and Q.~Berger,
\emph{Pinning of a renewal on a quenched renewal},
arXiv:1608.03265v2, 26 April 2017.
\url{https://arxiv.org/abs/1608.03265v2}.

\bibitem{Berger}
Q.~Berger,
\emph{Random Interfaces and Pinning Models}, lecture notes, 4 March 2026.
\url{https://www.math.univ-paris13.fr/~quentin.berger/documents/Random-Interfaces.pdf}.

\bibitem{JL}
S.~Junk and H.~Lacoin,
\emph{Coincidence of critical points for directed polymers for general
environments and random walks}, arXiv:2502.04113v2, 2025.
\url{https://arxiv.org/abs/2502.04113v2}.

\bibitem{Lacoin}
H.~Lacoin,
\emph{The martingale approach to disorder irrelevance for pinning models},
arXiv:1002.4753v2, 14 March 2011.
\url{https://arxiv.org/abs/1002.4753v2}.

\bibitem{Toninelli}
F.~L. Toninelli,
\emph{Disordered pinning models and copolymers: beyond annealed bounds},
Annals of Applied Probability \textbf{18} (2008), 1569--1587.
\url{https://doi.org/10.1214/07-AAP496}.
\end{thebibliography}
\end{document}
